% ===== 原文 pp.172–183（PDF pp.179–190） =====
% [原文 p.172]
\chapter{$A_\infty$ 范畴与三角范畴}

本章在引入 $A_\infty$ 代数、$A_\infty$ 范畴之后，讨论由 $A_\infty$ 范畴构造三角范畴的方法。$A_\infty$ 范畴以 DG 范畴为特例；当 $A_\infty$ 范畴是 DG 范畴时，由 $A_\infty$ 范畴构造三角范畴的方法就归结为 5.5 节中由 DG 范畴构造三角范畴的方法。如此，我们讨论由 $A_\infty$ 范畴得到的三角范畴所具有的各种性质。

\section{$A_\infty$ 代数}

以下固定域 $K$，$K$ 上的向量空间就简称为向量空间。为免繁琐，可以在全部讨论中都认为 $K$ 的示性数为零，特别地取 $K=\Q,\R,\C$ 等也无妨。实际上，我们所关心的对象首先就处于这种情形。

\subsection{$A_\infty$ 代数}

向分次向量空间 $V$ 中加入结合性结构而成的东西，称为（$K$ 上的）分次（结合）代数。这是分次微分（结合）代数（\textbf{D}ifferential \textbf{G}raded 代数，即 DG 代数）的微分为零的情形。而 DG 代数，则是（一般含有高阶积的）$A_\infty$ 代数在高阶积为平凡时的情形。

% [原文 p.173]
\begin{definition}{7.1}
$\Z$ 分次向量空间（graded vector space）是指以 $\Z$ 标号的向量空间 $V^{r}$（$r\in\Z$）的直和 $\bigoplus_{r\in\Z}V^{r}$。
\end{definition}

% [Tutor Note: 原文 p.172「$\Z$ ラベル付けられた」即「以 Z 标号/分次」；此处按定义直译，注意直和是「外直和」，$V^r$ 通过自然单射 $\iota_r$ 嵌入 $V$，这与定義 7.2 的同一视约定配套。]

以下，$\Z$ 分次向量空间就简称为分次向量空间。

例如，给定向量空间的复形 $\{V^{i},d^{i}\}_{i\in\Z}$ 时，$V:=\bigoplus_{i\in\Z}V^{i}$ 是分次向量空间。该复形的上同调的直和 $\bigoplus_{i\in\Z}H^{i}(V)$ 也是分次向量空间。

\begin{definition}{7.2}
设 $V=\bigoplus_{i\in\Z}V^{i}$ 为分次向量空间。对每个 $j\in\Z$，通过自然单射 $\iota_{j}\colon V^{j}\to V$，把 $V^{j}$ 的元素与 $\iota_{j}(V^{j})\subset V$ 视为同一。若对元素 $v\in V$，存在某个整数 $j\in\Z$ 使得 $v\in V^{j}$，则称 $v$ 是（关于次数）齐次的（homogeneous），并称 $j$ 为 $v$ 的次数（degree）。关于次数齐次的元素 $v\in V$ 的次数记作 $|v|$。
\end{definition}

由定义，分次向量空间 $V$ 的元素是 $V$ 的齐次元素的有限和。

在通常的不带次数的向量空间（2.4.3 节）的情形中，$K$ 双线性映射、多重线性映射等，都是把分次向量空间视为忘记了分次结构的通常向量空间时，按 2.4.3 节定义的那些东西。

\begin{definition}{7.3}
设 $V_{1},\dots,V_{k}$（$k\ge 1$），$W$ 为分次向量空间。作为分次向量空间的 $k$ 重线性映射
\[
  f\colon V_{1}\times\cdots\times V_{k}\longrightarrow W
\]
指的是作为向量空间的 $k$ 重线性映射。若 $k$ 重线性映射 $f\colon V_{1}\times\cdots\times V_{k}\to W$ 满足：存在某个整数 $r\in\Z$，使得对任意的齐次元素 $v_{1}\in V_{1},\dots,v_{k}\in V_{k}$ 都有
\[
  |f(v_{1},\dots,v_{k})|=|v_{1}|+\cdots+|v_{k}|+r,
\]
则称 $f$ 具有次数（degree）$r$。特别地，次数为零的多重线性映射称为保次数的多重线性映射。
\end{definition}

% [Tutor Note: 定义 7.3 的公式，原书 p.173 印作 $|f(v_1,\dots,v_n)|=|v_1|+\cdots+|v_n|+r$，右边末项下标为 $n$；按定义（$k$ 重、自变量为 $v_1,\dots,v_k$）应为 $|v_k|$。此处按上下文译为 $|v_k|$，并保留原印法于此备查。]

$k$ 重线性映射在 $k=1$ 时就是通常的线性映射，$k=2$ 时特别称为双线性映射。

\begin{definition}{7.4}
分次代数（graded algebra）$(A,m)$ 是指分次向量空间 $A$ 与保次数的双线性映射 $m\colon A\times A\to A$ 所成之组，且对任意元素 $a,b,c\in A$ 满足
\begin{equation}
  m(m(a,b),c)=m(a,m(b,c)). \tag{7.1}
\end{equation}
$m$ 称为分次代数 $(A,m)$ 的积（product），性质 (7.1) 则称 $m$ 是结合的（associative）。
\end{definition}

分次代数可以通过忘掉次数而视为代数，但要注意：与定义 2.71 中代数的定义不同，这里并未把单位元 $1_{A}$ 的存在列入条件。分次代数 $(A,m)$ 的单位元 $1_{A}$，是指它作为代数的单位元中次数（齐次地）为零者。以下，凡是涉及「有单位元的代数」，都如此说明。

\begin{definition}{7.5}
$(A,d,m)$ 是 DG 代数（DG-algebra），是指 $(A,m)$ 构成分次代数，且 $(A,d):=\{A^{r},d^{r}=d\}_{r\in\Z}$ 构成复形，并且 $d$ 与 $m$ 满足以下两个条件：对任意的 $a,b\in A$，
\[
  d^{2}=0,\qquad d m(a,b)=m(d(a),b)+(-1)^{|a|}m(a,d(b)).
\]
（此条件称为 DG 代数中的莱布尼茨法则（Leibniz rule）。）
\end{definition}

% [Tutor Note: 定义 7.5 的两个条件原文未编号；「$d$ 与 $m$ 满足以下两个条件」中的 $d^2=0$ 与复形条件重复，是原文写法，照译。$(-1)^{|a|}$ 的 $|a|$ 须理解为 $a$ 的齐次次数；对非齐次元按线性延拓。]

% [原文 p.174]
\begin{definition}{7.6}
（$A_\infty$ 代数（Stasheff$^{[50]}$））\ 称 $(A,\mathbf m:=\{m_{k}\}_{k\ge 1})$ 为 $A_\infty$ 代数，是指：$A=\bigoplus_{r\in\Z}A^{r}$ 是 $\Z$ 分次向量空间，$\mathbf m:=\{m_{n}\}_{n\ge 1}$ 是次数为 $(2-n)$ 的多重线性映射
\[
  m_{n}\colon A^{\otimes n}\longrightarrow A
\]
所构成的族，且对每个 $n=1,2,\dots$ 与关于次数齐次的元素 $a_{1},\dots,a_{n}\in A$，满足关系式
\begin{equation}
  0=\sum_{k+l=n+1}\ \sum_{j=0}^{k-1}\pm\,m_{k}\bigl(a_{1},\dots,a_{j},\ m_{l}(a_{j+1},\dots,a_{j+l}),\ a_{j+l+1},\dots,a_{n}\bigr) \tag{7.2}
\end{equation}
其中符号 $\pm=(-1)^{(j+1)(l+1)+l(|a_{1}|+\cdots+|a_{l}|)}$。
\end{definition}

% [Tutor Note: 式 (7.2) 的符号，底本 p.174 印作「± = (j+1)(l+1)+l(|a_1|+···+|a_l|)」（漏掉了指数 $(-1)^{\cdots}$）；据作者正誤表（2014-05-20 版）p.174 (7.2) 直後，正确形式为
%   ± = (−1)^{(j+1)(l+1)+l(|a_1|+···+|a_l|)}，
% 此处按订正后译出。核对方法：取 DG 代数（$m_{\ge 3}=0$）验证 $n=2$ 时退化出 $d(x\cdot y)=d(x)\cdot y+(-1)^{|x|}x\cdot d(y)$。]

我们来看 $A_\infty$ 关系式 (7.2) 在 $n=1,2,3$ 时的情形：把 $m_{1}=d$、$m_{2}=\cdot$ 记下，则对 $x,y,z\in A$ 有
\begin{align*}
  n=1)&\quad d^{2}=0,\\
  n=2)&\quad d(x\cdot y)=d(x)\cdot y+(-1)^{|x|}x\cdot d(y),\\
  n=3)&\quad (x\cdot y)\cdot z-x\cdot(y\cdot z)=d(m_{3})(x,y,z).
\end{align*}
这些分别意味着
\begin{itemize}[leftmargin=2em,itemsep=1pt,topsep=2pt]
  \item $n=1)$\ \ $(A,d)$ 构成复形，
  \item $n=2)$\ \ 微分 $d$ 关于积 $\cdot$ 满足莱布尼茨法则，
  \item $n=3)$\ \ 积 $\cdot$ 在同伦的意义下是结合的。
\end{itemize}
其中 $d(m_{3}):=m_{1}m_{3}+m_{3}(m_{1}\otimes 1\otimes 1+1\otimes m_{1}\otimes 1+1\otimes 1\otimes m_{1})$，$m_{3}$ 是 $(x\cdot y)\cdot z$ 与 $x\cdot(y\cdot z)$ 之间的同伦，而 $m_{4},m_{5},\dots$ 则被视为更高阶的同伦。

% [Tutor Note: $n=3)$ 的等号右边 $d(m_3)(x,y,z)$ 是「$m_3$ 与 $m_1$ 的复合」的简写，即 $m_1m_3+m_3(m_1\otimes 1\otimes 1+1\otimes m_1\otimes 1+1\otimes 1\otimes m_1)$；原文在此处把 $m_3$ 的次数约定（$|m_3|=2-3=-1$）省略了，这正是 $m_3$ 能充当同伦的原因。可验证：设 $A$ 为常结合代数（$m_{\ge 3}=0$），则 $|m_3|=0$ 与结合律相容。]

特别地，当 $m_{3}=0$ 时，积 $\cdot$ 成为结合的。没有高阶积的 $A_\infty$ 代数 $(A,\mathbf m)$，即 $m_{3}=m_{4}=\cdots=0$，就是 DG 代数（定义 7.5）。忘记这个 DG 代数的微分，$(A,m_{2})$ 便构成分次结合代数（定义 7.1.1）。特别地，$m_{1}=m_{3}=m_{4}=\cdots=0$ 的 $A_\infty$ 代数是分次结合代数。不过要注意：一般说到 $A_\infty$ 代数时，并不假定单位元 $1_{A}$ 的存在。

% [Tutor Note: 「定义 7.1.1」是原书 p.174 的误排，应为「定义 7.4」（分次代数）；作者正誤表未收此条。此处照译并于此指出，以免读者误以为存在 7.1.1 号定义。]

\begin{definition}{7.7}
设 $(A,\mathbf m)$ 为 $A_\infty$ 代数。若 $A$ 有次数为零的元素 $1_{A}\in A^{0}$，使得对任意元素 $a,a_{1},a_{2},\dots,a_{n}\in A$ 满足
\begin{itemize}[leftmargin=1.6em,itemsep=1pt,topsep=1pt]
  \item $m_{2}(1_{A},a)=m_{2}(a,1_{A})=a$，
  \item $m_{n}(a_{1},\dots,1_{A},\dots,a_{n})=0$，\quad $n=1,3,4,\dots$，
\end{itemize}
则称 $(A,\mathbf m)$ 为有单位元的（unital）$A_\infty$ 代数。
\end{definition}

% [Tutor Note: 定义 7.7 的第一条即 $1_A$ 是 $m_2$ 的双边单位元；第二条中 $n=2$ 不含 $1_A$（故不与定义冲突），$n=1$ 时意味着 $m_1(1_A)=0$。又有单位元 $A_\infty$ 代数不必有 $m_1(1_A)=0$ 之外的额外限制。]

% [原文 p.175]
其次，$m_{1}$ 的次数为 $1$、$m_{2}$ 的次数为零（$m_{2}$ 保次数），这从「$m_{1}$ 是微分、$m_{2}$ 定义积结构」来看是自然的；而 $m_{n}$ 的次数为 $(2-n)$ 也是其自然推广。正因如此，$A_\infty$ 代数的定义式 (7.2) 中才包含乍看复杂的符号 $\pm$。这个符号，可以通过借助悬垂（suspension）对 $A_\infty$ 代数作另一种描述来理解。

\begin{definition}{7.8}
（悬垂）对 $\Z$ 分次向量空间 $A$，把其次数整体减少 $1$ 而得的 $\Z$ 分次向量空间记作 $A[1]$，并记其同构映射为
\[
  s\colon A\longrightarrow A[1],\qquad s(A^{r})=(A[1])^{r-1}.
\]
称 $s(A)=A[1]$ 为 $A$ 的悬垂（suspension）。
\end{definition}

此时，$A$ 上的 $A_\infty$ 结构 $m_{k}$ 诱导出 $A[1]$ 上的多重线性映射 $s(m_{k})$：
\[
  s(m_{k})\bigl(s(a_{1}),\dots,s(a_{k})\bigr):=(-1)^{\sum_{i=1}^{k}(k-i)|s(a_{i})|}\,s\,m_{k}(a_{1},\dots,a_{k}).
\]
由此，$s(m_{k})$ 的次数不论 $k=1,2,\dots$ 都为 $1$，而 $A_\infty$ 关系式 (7.2) 可改写为
\begin{equation}
\begin{aligned}
  0=\sum_{k+l=n+1}\ \sum_{j=0}^{k-1}&(-1)^{|s(a_{1})|+\cdots+|s(a_{j})|}\\
  &\times s(m_{k})\bigl(s(a_{1}),\dots,s(a_{j}),\ s(m_{l})(s(a_{j+1}),\dots,s(a_{j+l})),\ s(a_{j+l+1}),\dots,s(a_{n})\bigr)
\end{aligned}
\tag{7.3}
\end{equation}
（其中 $|s(a_{i})|=|a_{i}|-1$。）也就是说，这里出现的符号 $(-1)^{|s(a_{1})|+\cdots+|s(a_{j})|}$，可以理解为次数为 $1$ 的 $s(m_{k})$ 越过 $s(a_{1}),\dots,s(a_{j})$ 时产生的符号（称为 Koszul 符号（Koszul sign））。

以下对 $a\in A$，把 $s(a)\in A[1]$ 简记为 $a\in A[1]$；写成 $A_\infty$ 代数 $(A[1],\mathbf m)$ 时，$\mathbf m_{k}$ 指的是 $s(m_{k})$。

% [Tutor Note: 式 (7.3) 与式 (7.2) 的等价性是本书符号约定的关键：多次使用 $s$ 后，$m_n$ 的次数从 $2-n$ 变成 $1$，$A_\infty$ 关系式遂成为「$m$ 是 $T^c(A[1])$ 上的微分」这一单一条件（见 7.1.2）。验证办法：$n=2$、$k=l=1$ 时 (7.3) 只剩 $s(m_1)s(m_1)=0$ 一项，与 (7.2) 的 $n=1$ 情形一致；$n=3$ 时逐项比对 $(-1)^{|s(a_1)|}$ 的位置即可。]

\subsection{棒构成}

上述 (7.3) 的表述，由于 $m_{n}$ 的次数不论 $n$ 都统一为 $1$，就使得 $A_\infty$ 代数的另一种描述——所谓棒构成——成为可能。棒构成把 $A_\infty$ 代数描述为某个分次微分余代数。

\begin{definition}{7.9}
（余代数，余结合）设 $C$ 为（分次）向量空间。若 $C$ 上定义的余积 $\Delta\colon C\to C\otimes C$ 是余结合的（coassociative），即
\[
  (\Delta\otimes\id_{C})\Delta=(\id_{C}\otimes\Delta)\Delta,
\]
则称 $C$ 为余代数（coalgebra）。
\end{definition}

\begin{definition}{7.10}
（余微分）设 $C$ 为余代数。次数为 $1$ 的线性算子 $Q\colon C\to C$ 若满足
\[
  \Delta Q=(Q\otimes\id_{C})\Delta+(\id_{C}\otimes Q)\Delta,
\]
则称 $Q$ 为余微分（coderivation）。不过要注意：对 $x,y\in C$，运算时的符号按 $(\id_{C}\otimes Q)(x\otimes y)=(-1)^{|x|}\bigl(x\otimes Q(y)\bigr)$ 定义。
\end{definition}

% [原文 p.176]
\begin{remark}{7.11}
余结合与余微分的定义，分别不过就是以下图表可交换而已。
\[
\begin{tikzcd}[row sep=2.6em, column sep=2.6em]
  C \arrow[r,"\Delta"] \arrow[d,"\Delta"'] & C\otimes C \arrow[d,"\Delta\otimes\id_{C}"] \\
  C\otimes C \arrow[r,"\id_{C}\otimes\Delta"'] & C\otimes C\otimes C
\end{tikzcd}
\qquad
\begin{tikzcd}[row sep=2.6em, column sep=2.6em]
  C \arrow[r,"Q"] \arrow[d,"\Delta"'] & C \arrow[d,"\Delta"] \\
  C\otimes C \arrow[r,"\id_{C}\otimes Q+Q\otimes\id_{C}"'] & C\otimes C
\end{tikzcd}
\]
把这些箭头的方向全部反转、把余积 $\Delta$ 换成积 $m$，就得到代数中结合性与微分的定义，这一点要注意。实际上，$C$ 的双对偶向量空间 $C^{*}=\Hom_{K}(C,K)$ 上就由此得到了结合代数与微分的结构。
\end{remark}

这里，作为余代数 $C$，考虑域 $K$ 上分次向量空间 $V=\bigoplus_{r\in\Z}V^{r}$ 的分次张量余代数
\[
  T^{c}(V):=\prod_{k=0}^{\infty}\bigl(V^{\otimes k}\bigr),
\]
其中 $V^{\otimes 0}:=K$。余积 $\Delta\colon T^{c}(V)\to T^{c}(V)\otimes T^{c}(V)$ 通过对 $v_{1}\otimes\cdots\otimes v_{n}\in V^{\otimes n}$ 定义
\[
  \Delta(v_{1}\otimes\cdots\otimes v_{n}):=\bigoplus_{j=0}^{n}(v_{1}\otimes\cdots\otimes v_{j})\otimes(v_{j+1}\otimes\cdots\otimes v_{n})
\]
而得到。不过这里例如 $j=0$ 的项，读作 $1_{K}\otimes(v_{1}\otimes\cdots\otimes v_{n})$。

对次数为 $1$ 的多重线性映射 $q_{k}\colon V^{\otimes k}\to V$，$T^{c}(V)$ 上的余微分 $Q_{k}\colon T^{c}(V)\to T^{c}(V)$ 如下定义：对 $v_{1}\otimes\cdots\otimes v_{n}\in V^{\otimes n}$，当 $n\le k$ 时令
\[
  Q_{k}(v_{1}\otimes\cdots\otimes v_{n}):=\sum_{j=0}^{n}v_{1}\otimes\cdots\otimes v_{j}\otimes q_{k}(v_{j+1}\otimes\cdots\otimes v_{j+k})\otimes v_{j+k+1}\otimes\cdots\otimes v_{n},
\]
而当 $k>n$ 时令 $Q_{k}(v_{1}\otimes\cdots\otimes v_{n})=0$，则可确认这将定义出余微分。实际上，$T^{c}(V)$ 上任意余微分 $Q$ 都可用一集多重线性映射 $\{q_{k}\}_{k\ge 0}$ 表为 $Q=Q_{0}+Q_{1}+\cdots$（cf.\ [54]）。

% [Tutor Note: 上段两处按作者正誤表（2014-05-20 版）p.176 条订正：(i) 自变量应为「关于次数齐次的元素 $v_1,\dots,v_n\in V$」；(ii) $Q_k$ 的求和中每项须乘 Koszul 符号 $(-1)^{|v_1|+\cdots+|v_j|}$，即
%   Q_k(v_1\otimes\cdots\otimes v_n)=\sum_{j=0}^{n}(-1)^{|v_1|+\cdots+|v_j|}\,v_1\otimes\cdots\otimes v_j\otimes q_k(\cdots)\otimes v_{j+k+1}\otimes\cdots\otimes v_n .
% 底本 p.176 原印无此符号。]
% [Tutor Note: 底本此处印作「$n\le k$ のときは」而求和 $j=0,\dots,n$ 中 $q_k$ 需取到 $v_{j+k}$，当 $n<k$ 时越界；按定义应为「$n\ge k$ のとき」。作者正誤表未收此条，此处照译本并指出。]

对 $A_\infty$ 代数 $(A,\mathbf m)$，考虑张量余代数 $T^{c}(A[1])$。也就是考虑上式中令 $V=A[1]$ 的情形。$A_\infty$ 结构 $m_{k}$（$k=1,2,\dots$）各自定义 $T^{c}(A[1])$ 上的余微分，把它记作 $\mathbf m_{k}$。此时，$A_\infty$ 关系式 (7.2) 等价于条件
\[
  (\mathbf m)^{2}=0,\qquad \mathbf m:=\mathbf m_{1}+\mathbf m_{2}+\cdots
\]
称 $(T^{c}(A[1]),\mathbf m,\Delta)$ 为 $(A,\mathbf m)$ 的棒构成（bar construction）。

% [Tutor Note: 「$m$ 的次数齐一为 1」是棒构成能写成一个微分的关键：$\mathbf m=\sum_k \mathbf m_k$ 是 $T^c(A[1])$ 上的余微分（次数 1），$(\mathbf m)^2=0$ 正是 $A_\infty$ 关系式的整体形式。可验证：$n=1$ 项给出 $\mathbf m_1^2=0$；$n=2$ 项给出 $\mathbf m_1\mathbf m_2+\mathbf m_2(\mathbf m_1\otimes 1+1\otimes\mathbf m_1)=0$（即莱布尼茨法则）。]

一般地，张量余代数 $C$ 上的余微分 $Q$ 若满足 $Q^{2}=0$，则 $(C,Q,\Delta)$

% [原文 p.177]
称为微分余代数（differential coalgebra）。当 $C=T^{c}(V)$ 时，微分余代数的微分可以写成 $Q=Q_{0}+Q_{1}+\cdots$。因此，微分余代数 $(T^{c}(V),Q,\Delta)$ 中 $Q_{0}=0$ 的情形对应于 $A_\infty$ 代数的棒构成。与 $Q_{0}\ne 0$ 的情形相对应的，有所谓弱 $A_\infty$ 代数（weak $A_\infty$-algebra）（cf.\ [13], [32]）；此时由于 $Q_{0}=\mathbf m_{0}$ 的存在，一般 $(m_{1})^{2}=0$ 不再成立。

% [Tutor Note: 棒构成给出的是「$Q_0=0$ 的微分余代数」与「$A_\infty$ 代数」的等价（Koszul 对偶的另一半）。弱 $A_\infty$ 代数保留 $Q_0\neq0$，其结构映射含 $m_0$（常数项），故 $m_1^2\neq0$ 一般成立；这与「有曲率的 DG 代数」$d^2=\omega$ 类似。]

\subsection{$A_\infty$ 映射}

借助上述棒构成，可以自然地定义两个 $A_\infty$ 代数之间映射的概念。

\begin{definition}{7.12}
对余代数 $C,C'$，若保次数的线性映射 $\mathbf f\colon C\to C'$ 满足
\[
  \Delta\mathbf f=(\mathbf f\otimes\mathbf f)\Delta,
\]
则称 $\mathbf f$ 为余代数之间的同态映射（coalgebra homomorphism）。
\end{definition}

对两个 $\Z$ 分次向量空间 $V,V'$，其张量余代数之间的同态映射 $\mathbf f\colon T^{c}(V)\to T^{c}(V')$ 与保次数的多重线性映射族 $\{f_{i}\colon V^{\otimes i}\to V'\}_{i=0,1,\dots}$ 一一对应 [54]。特别地，对 $v_{1}\otimes\cdots\otimes v_{n}\in V^{\otimes n}$，可定义
\[
  \mathbf f(v_{1}\otimes\cdots\otimes v_{n}):=\sum_{i\ge 1}\ \sum_{j_{1}+\cdots+j_{i}=n}f_{j_{1}}(v_{1},\dots,v_{j_{1}})\otimes f_{j_{2}}(v_{j_{1}+1},\dots,v_{j_{1}+j_{2}})\otimes\cdots\otimes f_{j_{i}}(\dots,v_{n}).
\]
这里特别考虑 $\mathbf f$ 中 $f_{0}=0$ 者。

\begin{definition}{7.13}
（$A_\infty$ 映射）对两个 $A_\infty$ 代数 $(A,\mathbf m)$，$(A',\mathbf m')$，$A_\infty$ 映射（$A_\infty$-morphism）$\mathbf f\colon(A,\mathbf m)\to(A',\mathbf m')$ 是指：$\{f_{i}\colon(A[1])^{\otimes i}\to A'[1]\}_{i=1,2,\dots}$ 作为余代数之间的同态映射的提升 $\tilde{\mathbf f}\colon T^{c}(A[1])\to T^{c}(A'[1])$，且 $\tilde{\mathbf f}$ 是复形 $(T^{c}(A[1]),\mathbf m)$、$(T^{c}(A'[1]),\mathbf m')$ 之间的映射，即
\begin{equation}
  \mathbf m'\circ\tilde{\mathbf f}=\tilde{\mathbf f}\circ\mathbf m. \tag{7.4}
\end{equation}
\end{definition}

该关系式 (7.4) 等价于：对任意的 $n=1,2,\dots$、$a_{1},\dots,a_{n}\in A[1]$，以下等式成立：
\begin{align*}
  &\sum_{i\ge 1}\ \sum_{k_{1}+\cdots+k_{i}=n}m'_{i}(f_{k_{1}}\otimes\cdots\otimes f_{k_{i}})(a_{1},\dots,a_{n})\\
  &\qquad\qquad=\sum_{\substack{i+1+j=k\\ i+l+j=n}}f_{k}\bigl(\mathbf 1^{\otimes i}\otimes m_{l}\otimes\mathbf 1^{\otimes j}\bigr)(a_{1},\dots,a_{n}).
\end{align*}
特别地，由此条件式在 $n=1$ 时的情形
\[
  m'_{1}f_{1}=f_{1}m_{1},
\]
即知 $f_{1}\colon(A[1],m_{1})\to(A'[1],m'_{1})$ 是复形之间的映射。

% [Tutor Note: 定义 7.13 与关系式 (7.4) 的要点：$A_\infty$ 映射不是「保积的线性映射」，而是「棒构成之间的微分映射」；其分量 $f_1,f_2,\dots$ 满足一族无穷多等式，$n=1$ 时退化为链映射条件 $m'_1f_1=f_1m_1$，$n=2$ 时给出 $m_2$ 的「同伦意义下的乘性」。可验证：若所有 $f_{\ge2}=0$，则 (7.4) 只剩 $f_1$ 的链映射条件。]

% [原文 p.178]
\begin{definition}{7.14}
设 $(A,\mathbf m),(A',\mathbf m')$ 分别为有单位元的 $A_\infty$ 代数。若 $A_\infty$ 映射 $\mathbf f\colon(A,\mathbf m)\to(A',\mathbf m')$ 满足
\[
  f_{1}(1_{A})=1_{A'},\qquad f_{n}(\dots,1_{A},\dots)=0,\quad n\ge 2,
\]
则称 $\mathbf f$ 保持单位元。
\end{definition}

% [Tutor Note: 作者正誤表（2014-05-20 版）p.178 起注明：「単位元」应读作「恒等元」（以下「単位元」＝「恒等元」）。作者此举意在把 $1_A$（代数的单位元）与范畴的恒等射区分开。本译本统一作「单位元」，读者见到 $1_A$ 时按「$m_2$ 的单位元」理解即可。]

\begin{definition}{7.15}
$A_\infty$ 映射 $\mathbf f\colon(A,\mathbf m)\to(A',\mathbf m')$ 若 $f_{1}\colon(A[1],m_{1})\to(A'[1],m'_{1})$ 是复形之间的拟同构（定义 4.50），则称为 $A_\infty$ 拟同构（$A_\infty$-quasi-isomorphism）。特别地，若 $f_{1}$ 是复形之间的同构，则称 $\mathbf f$ 为 $A_\infty$ 同构（$A_\infty$-isomorphism）。
\end{definition}

以下，把由 $f_{k}\colon(A[1])^{\otimes k}\to A'[1]$ 诱导的多重线性映射用同一记号 $f_{k}\colon A^{\otimes k}\to A'$（次数 $(1-k)$）表示。

如上所述，$A_\infty$ 代数 $(A,\mathbf m)$ 通过棒构成，由分次微分余代数 $(T^{c}(A[1]),\mathbf m,\Delta)$ 来定义；但当 $A$ 为有限维时，取这个余代数的双对所得的代数便是非交换自由代数，可以视为 $\Z$ 分次「非交换」空间 $A$ 上的多项式环。此时，余微分 $\mathbf m$ 被解释为 $\Z$ 分次「非交换」空间 $A$ 上的向量场与解积。这样的几何直观，在理解 $A_\infty$ 代数、乃至同伦代数一般中的各种定义与定理时，往往非常有力（cf.\ [26]）。

% [Tutor Note: 「非交换空间上的向量场与解积（divergence）」即 Kontsevich–Soibelman 的「非交换形式几何」观点：棒构成 $T^c(A[1])$ 的双对是形式幂级数环，$\mathbf m$ 是形式向量场，$(\mathbf m)^2=0$ 即「向量场的平方为零」（可积性）。可验算：$A$ 为一维（$A=K$ 集中于次数 0）时，$T^c(A[1])$ 的对偶为 $K[[t]]$，$m_1$ 对应 $\partial_t$，$m_2$ 对应 $t\partial_t$ 型向量场的分量。]

\subsection{极小模型定理与分解定理}

\begin{definition}{7.16}
若 $A_\infty$ 代数 $(A,\mathbf m)$ 的 $m_{1}$ 为平凡（$m_{1}=0$），则称其为极小的（minimal）。
\end{definition}

作为同伦代数基本定理的极小模型定理，对 $A_\infty$ 代数而言是由 Kadeishvili [25] 给出的。

\begin{theorem}{7.17}
（极小模型定理）对任意 $A_\infty$ 代数 $(A,\mathbf m)$，存在复形 $(A,m_{1})$ 的上同调上的 $A_\infty$ 代数 $(H(A),\mathbf m')$，以及 $A_\infty$ 拟同构 $\mathbf f\colon(H(A),\mathbf m')\to(A,\mathbf m)$。
\end{theorem}

这里，上述定理保证了复形之间的拟同构 $f_{1}\colon(H(A),m'_{1})\to(A,m_{1})$ 的存在，所以关于 $H(A)$ 上的 $A_\infty$ 结构 $\mathbf m'$，$m'_{1}$ 是平凡的。也就是说 $(H(A),\mathbf m')$ 是极小 $A_\infty$ 代数。此时，$A_\infty$ 代数 $(H(A),\mathbf m')$ 称为 $(A,\mathbf m)$ 的极小模型（minimal model）。$(A,\mathbf m)$ 的极小模型在 $H(A)$ 上的 $A_\infty$ 同构意义下唯一确定 [25]（参见 [26]）。

极小模型也有具体的构成方法。首先，对 $A_\infty$ 代数 $(A,\mathbf m)$，固定复形 $(A[1],d=m_{1})$ 与 $(H(A[1]),0)$ 之间的同伦等价映射（cf.\ 补题 3.21）。也就是说，对 $A[1]$ 与 $H(A[1])$ 之间的单射 $\iota$、满射 $\pi$：
\begin{equation}
  H(A[1])\ \underset{\pi}{\overset{\iota}{\rightleftarrows}}\ A[1] \tag{7.5}
\end{equation}
令 $P:=\iota\circ\pi$，再取线性映射 $h\colon(A[1])^{r}\to(A[1])^{r-1}$ 使得

% [Tutor Note: 极小模型定理（Kadeishvili）是「同伦传递定理」的一种：$H(A)$ 上存在 $A_\infty$ 结构使 $A$ 与 $H(A)$ 拟同构。构造用的数据 $(\iota,\pi,h)$ 即标准的「同伦等价数据」：$\pi\iota=\id$，$d h+h d=\id-\iota\pi$，$h^2=0$，$\pi h=0$，$h\iota=0$。补题 3.21 即给出这组数据的存在性。]

% [原文 p.179]
\[
  dh+hd=\id_{A[1]}-P
\]
这样的映射存在。由定义有 $P^{2}=P$。又，在两边作用 $d$ 可得 $d=dhd$，由此得 $(dh)^{2}=dh$、$(hd)^{2}=hd$。直和分解 $A[1]=P(A[1])\oplus dh(A[1])\oplus hd(A[1])$ 有时称为复形 $(A[1],d:=m_{1})$ 的霍奇分解（Hodge decomposition）。$P(A[1])\subset A[1]$ 给出 $H(A[1])$ 的代表元。

此时，首先归纳地如下定义 $A_\infty$ 拟同构 $\mathbf f\colon(H(A[1]),\mathbf m')\to(A[1],\mathbf m)$。令 $f_{1}=\iota\colon H(A[1])\to A[1]$，并对 $k\ge 2$ 令
\[
  f_{k}=-h\sum_{i\ge 2}\ \sum_{k_{1}+\cdots+k_{i}=k}m_{i}(f_{k_{1}}\otimes\cdots\otimes f_{k_{i}}).
\]
用它归纳地定义 $m'_{k}$（$k\ge 2$）：
\[
  m'_{k}=\pi\circ\sum_{i\ge 2}\ \sum_{k_{1}+\cdots+k_{i}=k}m_{i}(f_{k_{1}}\otimes\cdots\otimes f_{k_{i}}).
\]
注意：只要 $f_{2},\dots,f_{k-1}$ 已确定，$m'_{k}$ 就能定义。

取 $\alpha',\beta',\gamma',\dots\in H(A[1])$，具体地首先有
\[
  m'_{2}(\alpha',\beta')=\pi m_{2}\bigl(\iota(\alpha'),\iota(\beta')\bigr),
\]
即 $m'_{2}$ 是由 $m_{2}$ 自然诱导出的上同调上的积。其次，简记 $\alpha:=\iota(\alpha')$ 等，则由 $f_{2}(\alpha',\beta')=-h m_{2}(\alpha,\beta)$ 得
\begin{align*}
  m'_{3}(\alpha',\beta',\gamma')&=\pi m_{3}(\alpha,\beta,\gamma)+\pi m_{2}\bigl(f_{2}(\alpha',\beta'),\gamma\bigr)+\pi m_{2}\bigl(\alpha,f_{2}(\beta',\gamma')\bigr)\\
  &=\pi m_{3}(\alpha,\beta,\gamma)-\pi m_{2}\bigl(h m_{2}(\alpha,\beta),\gamma\bigr)-\pi m_{2}\bigl(\alpha,h m_{2}(\beta,\gamma)\bigr)
\end{align*}
同样可得到 $m'_{4},m'_{5},\dots$；它们有基于树图的记述，那种记述更便于直观把握（cf.\ [26]）。可以证明：如此构成的东西确实给出 $A_\infty$ 代数 $(H(A),\mathbf m')$ 与 $A_\infty$ 映射 $\mathbf f\colon(H(A),\mathbf m')\to(A,\mathbf m)$，这就给出了极小模型定理（定理 7.17）的（另一个）证明。又，通过考察这个极小模型的具体构成法，也可知道极小模型定理对有单位元的 $A_\infty$ 代数同样成立。

\begin{theorem}{7.18}
对任意有单位元的 $A_\infty$ 代数 $(A,\mathbf m)$，存在其上同调上的有单位元的 $A_\infty$ 代数 $(H(A),\mathbf m')$，以及保持单位元的 $A_\infty$ 拟同构 $\mathbf f\colon(H(A),\mathbf m')\to(A,\mathbf m)$。
\end{theorem}

% [Tutor Note: 上式 $m'_3$ 的第二、三项里的符号不影响结论（$\pm$ 由 Koszul 符号决定），但阶数不能乱：$m'_3$ 只能是 $m_3$ 与「$m_2$ 复合 $h$ 再复合 $m_2$」两类项，这由 $|h|=-1$、$|m_2|=0$、$|m_3|=-1$ 的次数计算唯一确定。可验证：$H(A)$ 上 $m'_2$ 结合，是因为 $m'_1=0$ 时关系式 (7.2) 的 $n=3$ 项只剩 $m'_2$ 的结合律偏差，而它等于 $0$（$h$ 的构造保证）。]

要注意，这些高阶项作为特别情形包含了被称为（高阶）Massey 积的东西（cf.\ [42]）。由于 $m'_{1}$ 平凡，由关于 $\mathbf m'$ 的 $A_\infty$ 关系式 (7.2) 可知 $m'_{2}$ 在 $H(A)$ 上是结合的。也就是说，$m_{2}$ 在上同调 $H(A)$ 上诱导出结合的积结构。如此，$A_\infty$ 代数 $(A,\mathbf m)$ 的上同调上所诱导的 $\Z$ 分次结合代数通常记作 $H(A,\mathbf m):=(H(A),H(m_{2}))$，其中 $H(m_{2}):=m'_{2}$。

% [原文 p.180]
\begin{definition}{7.19}
若 $A_\infty$ 代数 $(A,\mathbf m)$ 具有与 $H(A,\mathbf m)$ 同构的极小模型，则称 $(A,\mathbf m)$ 为形式的（formal）。
\end{definition}

也就是说，$A_\infty$ 代数是形式的，当且仅当存在 $A_\infty$ 拟同构 $H(A,\mathbf m)\to(A,\mathbf m)$。

\begin{remark}{7.20}
「极小」「形式的」这类名称来源于有理同伦论（cf.\ [13], [42], [55]）。
\end{remark}

以上极小模型的具体构成，是所谓同伦代数的同调扰动理论的一部分（[26] 及其参考文献）。同样包含极小模型定理、且更强有力的，有 $A_\infty$ 代数的分解定理。要注意：$K$ 向量空间的复形 $C:=\{C^{r},d^{r}\}_{r\in\Z}$ 可以看作 $m_{1}=d^{r}\colon C^{r}\to C^{r+1}$、$m_{2}=m_{3}=\cdots=0$ 的（没有单位元的）$A_\infty$ 代数。特别地，当该复形的上同调平凡（对所有的 $r\in\Z$ 有 $H^{r}(C)=0$）时，称 $C$ 为线性可缩（linear contractible）的 $A_\infty$ 代数。

\begin{theorem}{7.21}
（分解定理）体 $K$ 上任意的 $A_\infty$ 代数 $(A,\mathbf m)$，都与某个极小 $A_\infty$ 代数 $M$ 与线性可缩 $A_\infty$ 代数 $C$ 的直和 $M\oplus C$ 作为 $A_\infty$ 代数同构。
\end{theorem}

\begin{proof}
（概略）取复形 $(A,m_{1})$ 的一个霍奇分解 (7.5)。也就是说，固定一个作为分次向量空间的直和分解 $A=M\oplus C$，其中 $M:=PA$，$C:=(m_{1}h+hm_{1})A$。假设 $(A,\mathbf m)$ 的 $A_\infty$ 结构 $m_{2},\dots,m_{k}$ 已是 $M$ 上的多重线性映射。对保次数的映射 $f_{k+1}\colon(A[1])^{\otimes k}\to A[1]$，令 $\mathbf f:=\{f_{1}=\id_{A[1]},f_{2}=0,\dots,f_{k}=0,f_{k+1},f_{k+2}=0,\dots\}$，则 $\tilde{\mathbf f}\colon T^{c}(A[1])\to T^{c}(A[1])$ 是余代数的同构映射，且 $\mathbf f\colon(A,\mathbf m)\to(A,\mathbf m')$（其中 $\mathbf m':=\tilde{\mathbf f}\circ\mathbf m\circ\tilde{\mathbf f}^{-1}$）是 $A_\infty$ 同构映射。特别地 $m'_{2}=m_{2},\dots,m'_{k}=m_{k}$，并且可以证明能选取 $f_{k+1}$ 使 $m'_{k+1}$ 成为 $M$ 上的多重线性映射。如此以归纳论证完成证明（cf.\ [26], [31]）。\hfill$\square$
\end{proof}

作为这个分解定理的推论，不仅极小模型定理（定理 7.17）成立，还可以说：若存在 $A_\infty$ 拟同构 $\mathbf f\colon(A,\mathbf m)\to(A',\mathbf m')$，则存在反向的 $A_\infty$ 拟同构 $\mathbf g\colon(A',\mathbf m')\to(A,\mathbf m)$。由于 $A_\infty$ 拟同构的合成仍是 $A_\infty$ 拟同构，也就是说，$A_\infty$ 拟同构定义了 $A_\infty$ 代数之间的等价关系。

\subsection{$A_\infty$ 同伦等价}

我们定义 $A_\infty$ 代数之间的同伦等价。与复形的同伦等价等定义的情形一样，先定义 $A_\infty$ 映射之间的同伦，再用它定义 $A_\infty$ 代数之间的同伦等价。

DG 代数 $K[t,dt]:=(K[t,dt],d,m)$ 如下定义：令 $(K[t,dt],m)$ 为以 $t,dt$ 为变量的 $K$ 上多项式环模去理想 $K[t](dt)^{2}$ 所得的 $K$ 代数，其次数取为 $|t|=0$、$|dt|=1$，微分取为 $d(t):=dt$、$d(dt)=0$，则

% [原文 p.181]
$K[t,dt]$ 构成 DG 代数。（像这样积可交换的 DG 代数，称为 DG 交换代数。）

\begin{definition}{7.22}
设 $(A,\mathbf m)$ 为 $K$ 上的 $A_\infty$ 代数。如下定义 $A_\infty$ 代数 $(\tilde A,\tilde{\mathbf m})=(A,\mathbf m)\otimes_{K}K[t,dt]$。
\begin{itemize}[leftmargin=1.6em,itemsep=2pt,topsep=2pt]
  \item $\tilde A[1]:=A[1]\otimes_{K}K[t,dt]=(A\otimes_{K}K[t])\oplus(A\otimes_{K}K[t]dt)$。即 $\tilde A[1]$ 的元素是形如 $\tilde a=a\cdot b$（$a\in A[1]$，$b\in K[t,dt]$）者的有限和。次数由 $|\tilde a|:=|a|+|b|$ 定义。
  \item $\tilde A[1]$ 上的 $A_\infty$ 结构 $\{\tilde m_{k}\}_{k\ge 1}$ 定义为
  \begin{align*}
    k=1)\quad &\tilde m_{1}(a\cdot b):=m_{1}(a)\cdot b+(-1)^{|a|}a\cdot db,\\
    k\ge 2)\quad &\tilde m_{k}(a_{1}\cdot b_{1},\dots,a_{k}\cdot b_{k}):={}\\
    &\qquad (-1)^{\sum_{i=1}^{k}|b_{i}|(|a_{i+1}|+\cdots+|a_{k}|)}\,m_{k}(a_{1},\dots,a_{k})\cdot(b_{1}\cdots b_{k}).
  \end{align*}
\end{itemize}
\end{definition}

可以直接验证 $(\tilde A,\tilde{\mathbf m})$ 仍构成 $A_\infty$ 代数。

\begin{remark}{7.23}
这一构成把 $K[t,dt]$ 换成一般的分次微分交换代数也同样可以进行。
\end{remark}

$A_\infty$ 代数 $(\tilde A[1],\tilde{\mathbf m})$ 的元素 $\tilde a\in\tilde A[1]$ 可以表成不含 $dt$ 的项与含一个 $dt$ 的项之和 $a_{0}(t)+a_{1}(t)dt$。（由 $(dt)^{2}=0$，$dt$ 不可能出现两次以上。）对不含 $dt$ 的项 $a_{0}(t)$ 代入 $t=0$、$t=1$ 所得者，给出 $A[1]$ 的元素。把它们分别记作 $\tilde a|_{t=0}$、$\tilde a|_{t=1}$。

同样地，对 $A_\infty$ 代数 $(A,\mathbf m)$、$(A',\mathbf m')$，当给定 $A_\infty$ 写像 $\tilde{\mathbf f}\colon(\tilde A,\tilde{\mathbf m})\to(\tilde A',\tilde{\mathbf m}')$ 时，对 $a_{1},\dots,a_{k}\in A[1]$，$\tilde f_{k}(a_{1},\dots,a_{k})\in\tilde A'[1]$ 成为不含 $dt$ 的项 $f_{k}(t)(a_{1},\dots,a_{k})$ 与含一个 $dt$ 的项 $h_{k}(t)(a_{1},\dots,a_{k})\cdot dt$ 之和。此时 $\tilde{\mathbf f}|_{t=0}:=\{f_{k}(t=0)\}_{k\ge 1}$、$\tilde{\mathbf f}|_{t=1}:=\{f_{k}(t=1)\}_{k\ge 1}$ 定义从 $(A,\mathbf m)$ 到 $(A',\mathbf m')$ 的 $A_\infty$ 写像。实际上，在 $\tilde{\mathbf f}\colon T^{c}(A[1])\to T^{c}(\tilde A'[1])$ 中，把不含 $h_{k}(t)$ 的部分记作 $\mathbf f(t)$、含一个的部分记作 $\mathbf h(t)$，则由 $\tilde{\mathbf f}$ 是 $A_\infty$ 写像的条件 $\tilde{\mathbf m}'\circ\tilde{\mathbf f}=\tilde{\mathbf f}\circ\tilde{\mathbf m}$ 可得
\begin{align}
  \mathbf m'\circ\mathbf f(t)&=\mathbf f(t)\circ\mathbf m, \tag{7.6}\\
  \frac{d}{dt}\mathbf f&=-\bigl(\mathbf m'\circ\mathbf h(t)+\mathbf h(t)\circ\mathbf m\bigr). \tag{7.7}
\end{align}
由 (7.6) 可知 $\tilde{\mathbf f}|_{t=0}=\mathbf f(t=0)$、$\tilde{\mathbf f}|_{t=1}=\mathbf f(t=1)$ 是 $A_\infty$ 写像。条件 (7.6)、(7.7) 是 $\tilde{\mathbf f}$ 为 $A_\infty$ 写像的充要条件。

% [Tutor Note: 式 (7.6) 按作者正誤表（2014-05-20 版）p.181 条订正：底本原印作 $\tilde{\mathbf m}'\circ f(t)=f\circ m(t)$，应为 $\mathbf m'\circ\mathbf f(t)=\mathbf f(t)\circ\mathbf m$（去掉 $\tilde{\ }$，并把 $m(t)$ 写作 $\mathbf m$）。本译本已按订正后译出。]

\begin{definition}{7.24}
设 $(A,\mathbf m)$、$(A',\mathbf m')$ 为 $A_\infty$ 代数。若 $A_\infty$ 写像 $\mathbf f,\mathbf g\colon(A,\mathbf m)\to(A',\mathbf m')$ 满足：存在正整数 $n$ 以及从 $(A,\mathbf m)$ 到 $(A',\mathbf m')$ 的 $A_\infty$ 写像的列 $\mathbf f=\mathbf f^{(0)},\mathbf f^{(1)},\dots,\mathbf f^{(n-1)},\mathbf f^{(n)}=\mathbf g$，并且对每个 $i=0,\dots,n-1$ 存在 $A_\infty$ 写像 $\tilde{\mathbf f}^{(i,i+1)}\colon(\tilde A,\tilde{\mathbf m})\to(\tilde A',\tilde{\mathbf m}')$ 使得 $(\tilde{\mathbf f}^{(i,i+1)})|_{t=0}=\mathbf f^{(i)}$、$(\tilde{\mathbf f}^{(i,i+1)})|_{t=1}=\mathbf f^{(i+1)}$，则称 $\mathbf f$ 与 $\mathbf g$ 是 $A_\infty$ 同伦的（$A_\infty$-homotopic）。此时记作 $\mathbf f\overset{\mathrm{h.e.}}{\sim}\mathbf g$。
\end{definition}

% [原文 p.182]
\begin{definition}{7.25}
给定 $A_\infty$ 代数 $(A,\mathbf m)$ 与 $(A',\mathbf m')$。若存在 $A_\infty$ 写像 $\mathbf f\colon(A,\mathbf m)\to(A',\mathbf m')$、$\mathbf g\colon(A',\mathbf m')\to(A,\mathbf m)$ 使得
\[
  \mathbf g\circ\mathbf f\overset{\mathrm{h.e.}}{\sim}\id_{(A,\mathbf m)},\qquad
  \mathbf f\circ\mathbf g\overset{\mathrm{h.e.}}{\sim}\id_{(A',\mathbf m')},
\]
则称 $(A,\mathbf m)$ 与 $(A',\mathbf m')$ 是 $A_\infty$ 同伦等价的（$A_\infty$-homotopy equivalent）。此时称 $\mathbf f,\mathbf g$ 为 $A_\infty$ 同伦等价写像（$A_\infty$-homotopy equivalence）。
\end{definition}

\begin{theorem}{7.26}
设 $K$ 为特征零的域，$(A,\mathbf m)$、$(A',\mathbf m')$ 为 $K$ 上的 $A_\infty$ 代数。对它们之间的 $A_\infty$ 写像 $\mathbf f\colon(A,\mathbf m)\to(A',\mathbf m')$，以下二者等价：
\begin{itemize}[leftmargin=1.6em,itemsep=1pt,topsep=1pt]
  \item $\mathbf f$ 是 $A_\infty$ 拟同构写像；
  \item $\mathbf f$ 是 $A_\infty$ 同伦等价写像。\hfill$\square$
\end{itemize}
\end{theorem}

\begin{proof}
（概略）对 $(A,\mathbf m)$ 应用分解定理（定理 7.21），得 $A_\infty$ 同型 $\Phi\colon(A,\mathbf m)\simeq M\oplus C$。此时 $M\rightleftarrows M\oplus C$ 两者都给出 $A_\infty$ 拟同构写像，而 $\iota\circ\pi$ 与 $\id_{M\oplus C}$ 之间的 $A_\infty$ 同伦可以比较容易地构成。对 $(A',\mathbf m')$ 同样得 $A_\infty$ 同型 $\Psi\colon(A',\mathbf m')\simeq M'\oplus C'$ 与 $A_\infty$ 拟同构写像 $M'\rightleftarrows M'\oplus C'$。此时，若 $\mathbf f$ 是 $A_\infty$ 拟同构写像，则作为复形情形之讨论（命题 3.22）的自然推广，可证 $\mathbf f$ 是同伦等价写像。反之，给定 $A_\infty$ 同伦等价写像 $\mathbf f\colon(A,\mathbf m)\to(A',\mathbf m')$、$\mathbf f'\colon(A',\mathbf m')\to(A,\mathbf m)$ 时，可证它们给出 $M$ 与 $M'$ 之间的 $A_\infty$ 同型写像（cf.\ [31]）。\hfill$\square$
\end{proof}

\subsection{$A_\infty$ Maurer--Cartan 方程与模空间}

为定义 $A_\infty$ Maurer--Cartan 方程，需要把 $A_\infty$ 代数考虑为具有良好性质者。本章为简单起见，对任意 $K$ 上 $A_\infty$ 代数 $(A,\mathbf m)$，固定正整数 $n$，考虑可换 $K$ 代数 $K[h]/(h^{n})$（cf.\ 例 2.27）的极大理想 $\mathfrak m=h\cdot(K[h]/(h^{n}))$ 与 $A$ 的张量积 $A\otimes_{K}\mathfrak m$ 所成的 $A_\infty$ 代数。本章中单说 $A_\infty$ 代数 $(A,\mathbf m)$ 时，指的就是这样的 $A_\infty$ 代数。

% [Tutor Note: 这里 $A\otimes_K\mathfrak m$ 的作用是把 $A_\infty$ 结构「形式地」截断到 $h^n=0$，从而 $m_k(a,\dots,a)$ 在 $k\ge n$ 时自动为零，$m_*(a)$ 的无穷和才有意义（见注意 7.28）。这与形式模空间/形式形变的处理一致：$h$ 是形变参数，$h^2=0$ 时即普通 Maurer--Cartan 线性化。]

\begin{definition}{7.27}
对 $A_\infty$ 代数 $(A,\mathbf m)$，称
\[
  m_{*}(a):=m_{1}(a)+m_{2}(a,a)+m_{3}(a,a,a)+\cdots=0,\qquad a\in(A[1])^{0}
\]
为 $A_\infty$ Maurer--Cartan 方程（$A_\infty$-Maurer--Cartan equation）。其解所成集合记作
\[
  \mathrm{MC}(A,\mathbf m):=\bigl\{a\in(A[1])^{0}\ \bigm|\ m_{*}(a)=0\bigr\}.
\]
\end{definition}

\begin{remark}{7.28}
$m_{*}(a)$ 表面上由无穷和表示，但由于现在 $\mathfrak m=h\cdot(K[h]/(h^{n}))$ 的影响，当 $k\ge n$ 时 $m_{k}(a,\dots,a)=0$，故 $A_\infty$ Maurer--Cartan 方程在整个 $(A[1])^{0}$ 上有定义。
\end{remark}

% [原文 p.183]
\begin{definition}{7.29}
设 $(A,\mathbf m)$ 为 $A_\infty$ 代数。若 $a,b\in\mathrm{MC}(A,\mathbf m)$ 满足：存在正整数 $n$ 与 $\mathrm{MC}(A,\mathbf m)$ 的元素列 $a=a^{(0)},a^{(1)},\dots,a^{(n)}=b$，且对每个 $i=0,\dots,n-1$ 存在满足
\[
  \tilde a^{(i,i+1)}|_{t=0}=a^{(i)},\qquad \tilde a^{(i,i+1)}|_{t=1}=a^{(i+1)}
\]
的元素 $\tilde a^{(i,i+1)}\in\mathrm{MC}(\tilde A[1],\tilde{\mathbf m})^{0}$，则称 $a,b$ 是规范等价的（gauge equivalent）。此时记作 $a\overset{\mathrm{g.e.}}{\sim}b$。
\end{definition}

\begin{definition}{7.30}
对 $A_\infty$ 代数 $(A,\mathbf m)$，称集合
\[
  \mathcal M(A,\mathbf m):=\mathrm{MC}(A,\mathbf m)\big/\overset{\mathrm{g.e.}}{\sim}
\]
为 $A_\infty$ Maurer--Cartan 方程的模空间（moduli）。
\end{definition}

\begin{proposition}{7.31}
对 $A_\infty$ 代数 $(A,\mathbf m)$、$(A',\mathbf m')$ 与 $A_\infty$ 写像 $\mathbf f\colon(A,\mathbf m)\to(A',\mathbf m')$，令
\[
  f(a):=\sum_{n=1}^{\infty}f_{n}(a,\dots,a),\qquad a\in(A[1])^{0}.
\]
此时若 $a\in\mathrm{MC}(A,\mathbf m)$，则 $f(a)\in\mathrm{MC}(A',\mathbf m')$。
\end{proposition}

\begin{proof}
把 $T^{c}(A[1])$ 的元素 $e^{a}$ 定义为
\[
  e^{a}:=1+a+a\otimes a+a\otimes a\otimes a+\cdots,
\]
则 $\mathbf m(e^{a})=e^{a}\otimes m_{*}(a)\otimes e^{a}$。另一方面，由于 $\tilde{\mathbf f}(e^{a})=e^{f(a)}$，若 $a\in\mathrm{MC}(A,\mathbf m)$ 则
\[
  e^{f(a)}\otimes m'_{*}\bigl(f(a)\bigr)\otimes e^{f(a)}=\mathbf m'\circ\tilde{\mathbf f}(e^{a})=\tilde{\mathbf f}\circ\mathbf m(e^{a})=0,
\]
于是得证。\hfill$\square$
\end{proof}

\begin{proposition}{7.32}
给定 $A_\infty$ 代数 $(A,\mathbf m)$、$(A',\mathbf m')$ 与 $A_\infty$ 写像 $\mathbf f\colon(A,\mathbf m)\to(A',\mathbf m')$。此时若两个元素 $a,b\in\mathrm{MC}(A,\mathbf m)$ 规范等价（$a\overset{\mathrm{g.e.}}{\sim}b$），则 $f(a)\overset{\mathrm{g.e.}}{\sim}f(b)$。
\end{proposition}

\begin{proof}
由 $\mathbf f$ 如下构成 $A_\infty$ 写像 $\tilde{\mathbf f}\colon(\tilde A,\tilde{\mathbf m})\to(\tilde A',\tilde{\mathbf m}')$：对 $k=1,2,\dots$ 令
\[
  \tilde f_{k}(a_{1}\cdot b_{1},\dots,a_{k}\cdot b_{k}):=(-1)^{\sum_{i=1}^{k}|b_{i}|(|a_{i+1}|+\cdots+|a_{k}|)}f_{k}(a_{1},\dots,a_{k})\cdot(b_{1}\cdots b_{k})
\]
（$a_{i}\cdot b_{i}\in A[1]\otimes K[t,dt]$，$i=1,\dots,k$）。由补题 7.31，若 $\tilde a\in\mathrm{MC}(\tilde A,\tilde{\mathbf m})$ 则 $\tilde{\mathbf f}(\tilde a)\in\mathrm{MC}(\tilde A',\tilde{\mathbf m}')$；特别地由 $\tilde{\mathbf f}$ 的定义有 $(\tilde{\mathbf f}(\tilde a))|_{t=0}=\mathbf f(\tilde a|_{t=0})$、$(\tilde{\mathbf f}(\tilde a))|_{t=1}=\mathbf f(\tilde a|_{t=1})$，于是得证。\hfill$\square$
\end{proof}

由此，$A_\infty$ 写像 $\mathbf f\colon(A,\mathbf m)\to(A',\mathbf m')$ 诱导出相应模空间之间的写像。把它记作
\[
  \mathcal M(\mathbf f)\colon\mathcal M(A,\mathbf m)\longrightarrow\mathcal M(A',\mathbf m'),
\]
则由定义以下事实是显然的。
